\section{Assembly of the upper estimate}
\label{sec:assembly}
We prove \eqref{eq:main}, leaving the sharpness statement to the next
section. Fix any of the maximum simplices $S$ specified in
Theorem~\ref{thm:main}, and normalize that same simplex as in
Section~\ref{sec:enclosure}. The affine map preserves vertex count,
maximality, $e(K)$, $E(K,S)$, and the correspondence of centroids.

Suppose first that $0\le e(K)\le H_d^{-1}$. Proposition~\ref{prop:enclosure}
provides an actual-facet-supported simplex $P$ with relative projection
error $\delta=H_de(K)\in[0,1]$. Since $K$ has at most $n+r$ vertices,
Corollary~\ref{cor:caps} gives one point of $K$ in each cap of coordinate
loss at most
\[
 \rho=\delta^{1/k},\qquad k=\min(r,d-1).
\]
Both conclusions of Proposition~\ref{prop:retention} concern the same
$K,S,\rho$. Taking their minimum yields
\[
 E(K,S)\le\min(64nr,16n^2)\rho
 =T_{d,r}H_d^{1/k}e(K)^{1/k}.
\]
At $e(K)=0$, this gives $\rho=0$ and $E(K,S)=0$ without division by
any vanishing parameter. At $e(K)=H_d^{-1}$ the strict enclosure gate
remains strict, as checked in Proposition~\ref{prop:enclosure}.

For $e(K)>H_d^{-1}$, we have $(H_de(K))^{1/k}>1$ and
$T_{d,r}\ge n$. Therefore the universal replacement bound $E(K,S)\le n$
proves the same estimate. Undoing the affine map yields \eqref{eq:main}
for the originally prescribed maximum and its original centroid.

The exponent is supplied by the finite apex count; the constants $64nr$
and $16n^2$ come from retention. Their minimum does not select different
maxima or different centers. No classification of all polytopes with two
extra vertices, stronger all-minor stochastic conjecture, or forest
conjecture enters this argument.
